\section{Conclusion}
\label{sec:conclusion}

We began with an 8.57\% worst-group accuracy --- a number that demands explanation. A DINO model trained on Waterbirds sees 97\% of waterbirds correctly when the background is water, and 8.57\% when the background is land. This is not a capacity failure. The 62\% average accuracy confirms that the model has learned; the 73.4 pp worst-group accuracy gap (73.36 pp exact) confirms that it has learned the wrong thing, with geometric reliability.

This work proposes and partially validates a geometric explanation: SSL pretraining with SGD creates directional curvature asymmetry in the InfoNCE loss landscape, privileging spurious feature directions over core feature directions. We call this sharpness anisotropy, and we design an annotation-free protocol to measure it --- using SAM perturbations to probe directional curvature and an LFR-style linear probe loss proxy to identify spurious directions without group labels.

Our confirmed contribution is twofold. First, we establish the SSL shortcut problem's severity and geometric stability: the 8.57\% WGA and 73.4 pp gap persist throughout training, ruling out convergence as a remedy. Second, we introduce the sharpness anisotropy ratio as a geometric diagnostic for SSL shortcut encoding, together with a SAM-SSL intervention protocol --- both designed and awaiting full experimental validation.

The mechanistic test (sharpness anisotropy measurement and SAM-SSL training) is ongoing, pending the full 200-epoch experimental runs. The theoretical prediction is clear: if anisotropy exists and SAM flattens it, WGA should improve. If anisotropy does not exist --- if the InfoNCE loss landscape does not exhibit directional curvature asymmetry --- we will have identified an important boundary condition of geometric shortcut theory, ruling out optimizer geometry as the mechanism and pointing toward representational or objective-level explanations.

\textbf{Future Directions.} Measuring anisotropy emergence across epochs (50, 100, 150, 200) will reveal whether spurious encoding is gradual or abrupt. Validating the LFR proxy precision/recall against Waterbirds group labels will determine whether annotation-free geometric diagnosis is reliable in SSL settings. The full 9-combination evaluation (3 SSL methods $\times$ 3 datasets) will establish whether the SSL shortcut problem is SSL-method-specific or broadly characteristic of InfoNCE-based pretraining on spurious data.

We hope this work encourages the community to view SSL spurious correlation robustness as a geometric problem --- not just a representational one --- and to explore optimizer-level interventions as a principled alternative to post-hoc debiasing on fixed representations. The geometry of learning, not just what is learned, determines who a model fails.
